% 彗星（综述）
% license CCBYSA3
% type Wiki

本文根据 CC-BY-SA 协议转载翻译自维基百科\href{https://en.wikipedia.org/wiki/Comet}{相关文章}。

\begin{figure}[ht]
\centering
\includegraphics[width=8cm]{./figures/45ca94c34375617b.png}
\caption{} \label{fig_Comet_1}
\end{figure}
彗星是一类由冰组成的、体积较小的太阳系天体或星际天体。当其经过接近太阳的区域时，会因受热而开始释放气体，这一过程称为释气作用（outgassing）。释气作用会在彗核周围形成一个延展的、未被引力束缚的大气层，即彗发（coma）；在某些情况下，彗发中的气体与尘埃会被吹离，形成彗尾（tail）。这些现象均由太阳辐射以及向外流动的太阳风等离子体作用于彗核所引起。

彗核的尺寸从数百米至数十公里不等，由松散聚集的冰、尘埃与微小岩石颗粒构成。彗发的尺度可达地球直径的 15 倍，而彗尾甚至可以延伸至超过 1 个天文单位。如果彗星距离地球足够近且亮度足够高，人类可以在无望远镜的情况下直接观测到它，其在天空中的角距可达 30 度（相当于 60 个满月）。自古以来，彗星已被许多文化与宗教所观察与记载。

彗星通常具有高离心率的椭圆轨道，其公转周期跨度极大，从数年到可能长达数百万年不等。短周期彗星起源于海王星轨道之外的柯伊伯带（Kuiper belt）或相关的散射盘（scattered disc）。长周期彗星被认为来自奥尔特云（Oort cloud），该云由冰质天体组成，其范围从柯伊伯带外侧一直延展至通往最近恒星方向一半的距离。长周期彗星通常因邻近恒星的引力扰动或银河潮汐效应而被触发向太阳内侧运动。

双曲轨道彗星（hyperbolic comets）可能仅一次穿越内太阳系，随后被抛入星际空间。彗星的出现被称为一次彗星显现（apparition）。
